\subsection{Annäherung statt Auswertung}
\label{subsec:real_functions_limit_intuition}
Wir beginnen mit der bekannten Funktion $f(x)=x^2$.
Die Eingaben nähern sich der Stelle $2$ von links und von rechts;
den Wert bei $x=2$ lassen wir zunächst aus.
\begin{center}
  \renewcommand{\arraystretch}{1.4}
  \begin{tabular}{@{}c|rrr@{\qquad}rrr@{}}
    $x$ & $1{,}9$ & $1{,}99$ & $1{,}999$ & $2{,}001$ & $2{,}01$ & $2{,}1$\\
    $x^2$ & $3{,}61$ & $3{,}9601$ & $3{,}996001$ &
      $4{,}004001$ & $4{,}0401$ & $4{,}41$\\
  \end{tabular}
\end{center}
\begin{center}
  \CachedTikzFigure{content/resources/chapters/real_functions/tikz/functions_limit_approach.tex}
\end{center}

\par\noindent\begin{minipage}{\linewidth}
\begin{DefinitionBox}{Grenzwert -- anschauliche Bedeutung}{real_functions_limit}
  Die Funktion $f$ hat für $x\to a$ den \emph{Grenzwert}
  \EnglishTerm{limit} $L\in\R$, wenn ihre Werte beliebig nahe bei $L$
  liegen und bleiben, sobald $x$ hinreichend nahe bei $a$ liegt,
  aber $x\ne a$ ist. Wir schreiben
  \[
    \lim_{x\to a}f(x)=L
    \qquad\text{oder}\qquad
    f(x)\to L\quad\text{für }x\to a.
  \]
  Die Annäherung erfolgt nur mit $x\in D_f$; solche Stellen müssen
  beliebig nahe bei $a$ liegen.
\end{DefinitionBox}
\begin{ExampleBox}[before skip=0.3\baselineskip]{Die Annäherung begründen}{real_functions_square_limit}
  Warum bleiben die Werte von $x^2$ nahe bei $4$, wenn $x$ nahe bei $2$ liegt?
  \begin{ClozeArea}[34]{6}
    Die Differenz lässt sich mit einer binomischen Formel zerlegen:
    \[
      |x^2-4|=|x-2|\,|x+2|.
    \]
    Für $1<x<3$ ist $|x+2|<5$. Der Abstand zu $4$ ist also kleiner als
    das Fünffache des Abstands von $x$ zu $2$. Wird letzterer beliebig
    klein, gilt dies auch für den ersten:
    $\lim_{x\to2}x^2=4$.
  \end{ClozeArea}
\end{ExampleBox}
\end{minipage}


\begin{KeyPointBox}{Der Wert an der Stelle ist nicht die Frage}{real_functions_limit_vs_value}
  \begin{center}
    \CachedTikzFigure{content/resources/chapters/real_functions/tikz/functions_limit_questions.tex}
  \end{center}
  Für den Grenzwert ist nicht entscheidend, ob $f(a)$ existiert oder
  welchen Wert es hat. Entscheidend ist das Verhalten in der Nähe von $a$.
\end{KeyPointBox}
\begin{RemarkBox}
  Eine Wertetabelle oder Zeichnung hilft beim Erkennen eines Grenzwerts.
  Endlich viele Stichproben beweisen ihn jedoch nicht: Die Aussage muss
  für jede hinreichend nahe zulässige Eingabe gelten.
\end{RemarkBox}
